% from ega4-17-fr.tex
\oldpage[IV-4]{61}
$\mathop{Gr}\nolimits_1(g)=0$ (17.2.1). En outre, $g$ est localement de présentation finie (I, 4.3.1); donc l'hypothèse $\mathop{Gr}\nolimits_1(g)=0$ entraîne que $g$ est une immersion ouverte (16.1.10); étant surjective, cette immersion est un isomorphisme, donc $f$ est un monomorphisme (I, 5.3.8).

\subsection{Morphismes lisses, morphismes non ramifiés, morphismes étales}
\EGAReaderTarget{ega4.c04.subsec.17-3}{subsection:IV.17.3}

\begin{definition}[17.3.1]
\EGAReaderTarget{ega4.c04.def.17-3-1}{IV.17.3.1}
On dit qu'un morphisme $f:X\to Y$ est \emph{lisse} (resp. \emph{non ramifié}, ou \emph{net}\footnote{Les mots « net » et « formellement net » paraissent bien préférables à la terminologie consacrée de « non ramifié » (resp. « formellement non ramifié ») et seront employés à peu près exclusivement à partir du chap. V. Dans ce chapitre, nous avons conservé l'ancienne terminologie afin de ne pas entrer en conflit avec (0, 19.10).}, resp. \emph{étale}) s'il est localement de présentation finie et formellement lisse (resp. formellement non ramifié, resp. formellement étale).
\end{definition}

On dit encore alors que $X$ est \emph{lisse} (resp. \emph{non ramifié} ou \emph{net}, resp. \emph{étale}) sur $Y$.

Nous verrons plus loin (17.5.2) que cette définition d'un morphisme lisse coïncide avec celle déjà donnée dans (6.8.1); jusque-là, c'est la définition de (17.3.1) que nous utiliserons exclusivement.

Il est clair que dire que $f$ est étale signifie qu'il est \emph{à la fois} lisse et non ramifié.

\begin{remarks}[17.3.2]
\EGAReaderTarget{ega4.c04.subsec.17-3-2}{IV.17.3.2}
\begin{enumerate}
  \item[(i)] On notera que la définition (17.3.1) peut s'exprimer uniquement à l'aide du foncteur
  \[
    Y'\longmapsto\Hom_Y(Y',X)
  \]
  considéré dans (17.1.2, (iii)), car dire que $f$ est localement de présentation finie équivaut à dire que le foncteur précédent \emph{commute aux limites projectives de schémas affines} (8.14.2).
  \item[(ii)] Soient $A$ un anneau, $B$ une $A$-algèbre. On dit que $B$ est une $A$-algèbre \emph{lisse} (resp. \emph{non ramifiée}, resp. \emph{étale}) si le morphisme correspondant $\Spec(B)\to\Spec(A)$ est lisse (resp. non ramifié, resp. étale). Il revient au même de dire que $B$ est une $A$-algèbre \emph{de présentation finie} (I, 4.4.6) et formellement lisse (resp. formellement non ramifiée, resp. formellement étale) pour les topologies discrètes.
  \item[(iii)] Il résulte de (17.1.6) et de la définition d'un morphisme localement de présentation finie (I, 4.4.2) que les notions de morphisme lisse, non ramifié et étale sont \emph{locales sur $X$ et sur $Y$}.
\end{enumerate}
\end{remarks}

\begin{proposition}[17.3.3]
\label{IV.17.3.3}
\begin{enumerate}
  \item[(i)] Une immersion ouverte est étale. Pour qu'une immersion soit non ramifiée, il faut et il suffit qu'elle soit localement de présentation finie.
  \item[(ii)] Le composé de deux morphismes lisses (resp. non ramifiés, resp. étales) est lisse (resp. non ramifié, resp. étale).
  \item[(iii)] Si $f:X\to Y$ est un $S$-morphisme lisse (resp. non ramifié, resp. étale), il en est de même de $f_{(S')}:X_{(S')}\to Y_{(S')}$ pour toute extension $S'\to S$ du préschéma de base.
  \item[(iv)] Si $f:X\to X'$ et $g:Y\to Y'$ sont deux $S$-morphismes lisses (resp. non ramifiés, resp. étales), il en est de même de $f\times_S g:X\times_S Y\to X'\times_S Y'$.
\end{enumerate}

